\documentclass[11pt,a4paper]{article}

\usepackage[utf8]{inputenc}
\usepackage[
  left=1.05in,
  right=1.05in,
  top=1in,
  bottom=1in,
  headheight=14pt
]{geometry}
\usepackage{amsmath,amssymb,amsthm,mathtools}
\usepackage{bm}
\usepackage{enumitem}
\usepackage{booktabs,array}
\usepackage{microtype}
\usepackage{xcolor}
\usepackage[hidelinks]{hyperref}
\usepackage{fancyhdr}

\linespread{1.10}
\setlength{\parindent}{0pt}
\setlength{\parskip}{0.48em plus 0.12em minus 0.08em}
\setlength{\abovedisplayskip}{0.8em plus 0.2em minus 0.15em}
\setlength{\belowdisplayskip}{0.8em plus 0.2em minus 0.15em}
\setlength{\abovedisplayshortskip}{0.55em plus 0.15em}
\setlength{\belowdisplayshortskip}{0.65em plus 0.15em}
\setlength{\jot}{5pt}
\setlength{\emergencystretch}{2em}
\raggedbottom

\numberwithin{equation}{section}

\newtheoremstyle{spacedplain}
  {8pt}{8pt}{\itshape}{}{\bfseries}{.}{0.5em}{}
\newtheoremstyle{spaceddefinition}
  {8pt}{8pt}{\normalfont}{}{\bfseries}{.}{0.5em}{}
\newtheoremstyle{spacedremark}
  {8pt}{8pt}{\normalfont}{}{\itshape}{.}{0.5em}{}

\theoremstyle{spacedplain}
\newtheorem{theorem}{Theorem}[section]
\newtheorem{lemma}[theorem]{Lemma}
\newtheorem{proposition}[theorem]{Proposition}
\newtheorem{corollary}[theorem]{Corollary}

\theoremstyle{spaceddefinition}
\newtheorem{definition}[theorem]{Definition}
\newtheorem{example}[theorem]{Example}
\newtheorem{problem}[theorem]{Problem}

\theoremstyle{spacedremark}
\newtheorem{remark}[theorem]{Remark}

\setlist[itemize]{leftmargin=*,itemsep=4pt,topsep=6pt,parsep=2pt}
\setlist[enumerate]{leftmargin=*,itemsep=4pt,topsep=6pt,parsep=2pt}

\pagestyle{fancy}
\fancyhf{}
\fancyhead[L]{MTL712}
\fancyhead[R]{Stability and Difference Equations}
\fancyfoot[C]{\thepage}
\renewcommand{\headrulewidth}{0.4pt}
\renewcommand{\footrulewidth}{0pt}

\title{\textcolor{blue!50!black}{\textbf{MTL712: Computational Methods for Differential\\Equations}}\\[6pt]
\large Stability, Difference Equations, and Order Conditions}
\author{\normalsize Cleaned transcription of handwritten notes}
\date{}

\begin{document}

\maketitle
\thispagestyle{plain}
\tableofcontents
\clearpage

\section{Root Condition and Zero-Stability}

Consider the \(k\)-step linear multistep method
\[
  \sum_{j=0}^{k}\alpha_j y_{n+j}
  =h\sum_{j=0}^{k}\beta_j f_{n+j},
  \qquad
  \alpha_k\neq 0,
\]
with first characteristic polynomial
\[
  \rho(z)=\sum_{j=0}^{k}\alpha_jz^j.
\]

\begin{proof}[Proof of the root condition]
Apply the method to \(y'=0\).  The
resulting recurrence is
\[
  \alpha_k y_{n+k}+\alpha_{k-1}y_{n+k-1}
  +\cdots+\alpha_1y_{n+1}+\alpha_0y_n=0.
\]
If \(z_s\) is a zero of \(\rho\) of multiplicity \(m_s\), the
corresponding terms in the general solution have the form
\[
  p_s(n)z_s^n,
  \qquad \deg p_s\leq m_s-1.
\]
If \(\lvert z_s\rvert>1\), some solutions therefore grow like
\(\lvert z_s\rvert^n\).  If \(\lvert z_s\rvert=1\) and \(m_s>1\), some
solutions grow like \(n^{m_s-1}\).  In either case they are unbounded as
\(n\to\infty\), even though \(nh=b-a\) is held fixed while \(h\to0\).

Choose initial values \(y_0,\ldots,y_{k-1}\) that generate such an
unbounded solution, and compare them with
\(\widehat y_0=\cdots=\widehat y_{k-1}=0\), whose corresponding solution
is \(\widehat y_n=0\).  No constant \(K\), independent of \(h\), can
satisfy
\[
  \lvert y_n-\widehat y_n\rvert
  \leq K\max_{0\leq j\leq k-1}
  \lvert y_j-\widehat y_j\rvert .
\]
Thus violation of the root condition implies failure of zero-stability.
For the converse, the root condition implies that the powers of the
companion matrix of the homogeneous recurrence are uniformly bounded:
roots strictly inside the unit disk decay, while simple roots on the
unit circle remain bounded.  Applying the variation-of-constants
formula to the difference of two numerical solutions, followed by the
Lipschitz condition and a discrete Gronwall inequality, gives the
required bound with a constant independent of \(h\).  Hence the method
is zero-stable.
\end{proof}

\begin{example}[Euler's method]
For
\[
  y_{n+1}=y_n+h f(x_n,y_n),
\]
the first characteristic polynomial is \(\rho(z)=z-1\).  Its only zero,
\(z=1\), is simple, so Euler's method is zero-stable.
\end{example}

\begin{example}[A non-zero-stable three-step method]
Consider the sixth-order formula
\[
  11y_{n+3}+27y_{n+2}-27y_{n+1}-11y_n
  =3h\bigl(f_{n+3}+9f_{n+2}+9f_{n+1}+f_n\bigr).
\]
Its first characteristic polynomial is
\[
  \rho(z)=11z^3+27z^2-27z-11
  =(z-1)(11z^2+38z+11).
\]
The roots are
\[
  z_1=1,\qquad
  z_{2,3}=\frac{-19\pm2\sqrt{60}}{11}
  \approx-0.3189,\,-3.1357.
\]
Since one root has modulus greater than \(1\), the method is not
zero-stable.
\end{example}

\section{Stability of Linear Multistep Methods}

Let \(E\) denote the shift operator, \(E y_n=y_{n+1}\), and let
\[
  \sigma(z)=\sum_{j=0}^{k}\beta_jz^j
\]
be the second characteristic polynomial.  The multistep method can be
written as
\[
  \rho(E)y_n=h\sigma(E)f_n,
  \qquad\text{or}\qquad
  \rho(E)y_n-h\sigma(E)f_n=0.
\]
Apply it to the test equation
\[
  y'=\lambda y,\qquad \lambda<0.
\]
Then
\[
  \bigl[\rho(E)-\bar h\,\sigma(E)\bigr]y_n=0,
  \qquad \bar h=h\lambda,
\]
whose stability polynomial is
\[
  \pi(z,\bar h)=\rho(z)-\bar h\,\sigma(z).
\]
Let \(z_1,\ldots,z_k\) be its zeros, numbered so that \(z_1\to1\) as
\(\bar h\to0\) for a consistent zero-stable method.

\begin{definition}[Absolute and relative stability]
At a given value of \(\bar h\), the method is \emph{absolutely stable}
if
\[
  |z_r|<1,\qquad r=1,\ldots,k.
\]
It is \emph{relatively stable} if
\[
  |z_r|<|z_1|,\qquad r=2,\ldots,k.
\]
\end{definition}

\begin{definition}[Interval of absolute stability]
An interval \((a,b)\) on the real axis is an interval of absolute
stability if the method is absolutely stable for every
\(\bar h\in(a,b)\).
\end{definition}

\begin{definition}[Strong and weak stability]
A zero-stable multistep method is \emph{strongly stable} if \(z=1\) is
the only zero of \(\rho\) on the unit circle.  It is \emph{weakly
stable} if \(\rho\) has additional (necessarily simple) zeros on the
unit circle.
\end{definition}

Two standard tools for locating an interval of absolute stability are
Schur's criterion and the Routh--Hurwitz criterion.

\subsection{Schur's Criterion}

Let
\[
  \phi(z)=c_0z^k+c_1z^{k-1}+\cdots+c_{k-1}z+c_k,
  \qquad c_0c_k\neq0,
\]
have complex coefficients, and define its conjugate-reciprocal
polynomial by
\[
  \widehat\phi(z)
  =\overline{c_k}z^k+\overline{c_{k-1}}z^{k-1}
   +\cdots+\overline{c_1}z+\overline{c_0}.
\]

\begin{definition}
\(\phi\) is a \emph{Schur polynomial} if each of its zeros \(z_r\)
satisfies \(|z_r|<1\).
\end{definition}

Define the polynomial of degree \(k-1\)
\[
  \phi_1(z)
  =\frac{\widehat\phi(0)\phi(z)-\phi(0)\widehat\phi(z)}{z}.
\]

\begin{theorem}[Schur's criterion]
The polynomial \(\phi\) is a Schur polynomial if and only if
\[
  |\widehat\phi(0)|>|\phi(0)|
  \quad\text{and}\quad
  \phi_1\ \text{is a Schur polynomial}.
\]
Equivalently, the first inequality is \(|c_0|>|c_k|\).
\end{theorem}

\subsection{Routh--Hurwitz Criterion}

The Cayley transformation
\[
  w=\frac{z-1}{z+1},
  \qquad
  z=\frac{1+w}{1-w},
\]
maps the open unit disk in the \(z\)-plane onto the open left
half-plane \(\operatorname{Re}w<0\).  Under this transformation the
stability polynomial \(\pi(z,\bar h)\) becomes, after multiplication by
\((1-w)^k\), a polynomial
\[
  q(w)=a_0w^k+a_1w^{k-1}+\cdots+a_k.
\]
Thus the zeros of \(\pi\) lie in \(|z|<1\) if and only if the zeros of
\(q\) lie in \(\operatorname{Re}w<0\).

\begin{theorem}[Routh--Hurwitz criterion]
Assume \(a_0>0\) and put \(a_j=0\) for \(j>k\).  All zeros of \(q\) lie
in the open left half-plane if and only if all leading principal minors
of the Hurwitz matrix
\[
  H=
  \begin{pmatrix}
    a_1&a_3&a_5&\cdots&a_{2k-1}\\
    a_0&a_2&a_4&\cdots&a_{2k-2}\\
    0&a_1&a_3&\cdots&a_{2k-3}\\
    0&a_0&a_2&\cdots&a_{2k-4}\\
    \vdots&\vdots&\vdots&\ddots&\vdots
  \end{pmatrix}_{k\times k}
\]
are positive.
\end{theorem}

In particular:
\begin{itemize}
  \item for \(k=2\): \(a_0,a_1,a_2>0\);
  \item for \(k=3\): \(a_0,a_1,a_2,a_3>0\) and
    \(a_1a_2-a_0a_3>0\);
  \item for \(k=4\): \(a_0,a_1,a_2,a_3,a_4>0\) and
    \[
      a_1a_2a_3-a_0a_3^2-a_1^2a_4>0.
    \]
\end{itemize}

\section{Difference Operators, Consistency, and Order}

Associated with the \(k\)-step method is the linear difference operator
\[
\begin{split}
  L[y(x);h]
  ={}&\sum_{j=0}^{k}\alpha_jy(x+jh)
  -h\sum_{j=0}^{k}\beta_jy'(x+jh).
\end{split}
\]
For a sufficiently differentiable function \(y\), Taylor expansion
gives
\[
  L[y(x);h]
  =c_0y(x)+c_1hy'(x)+c_2h^2y''(x)+\cdots,
\]
where
\[
  c_0=\sum_{j=0}^{k}\alpha_j,
  \qquad
  c_1=\sum_{j=0}^{k}j\alpha_j-\sum_{j=0}^{k}\beta_j,
\]
and, for \(q\geq2\),
\[
  c_q
  =\frac{1}{q!}\sum_{j=0}^{k}j^q\alpha_j
   -\frac{1}{(q-1)!}\sum_{j=0}^{k}j^{q-1}\beta_j.
\]
In terms of the characteristic polynomials,
\[
  c_0=\rho(1),
  \qquad
  c_1=\rho'(1)-\sigma(1).
\]

\begin{definition}[Order of a difference operator]
The difference operator \(L\) has order \(p\) if
\[
  c_0=c_1=\cdots=c_p=0,
  \qquad c_{p+1}\neq0.
\]
\end{definition}

The same method may be written with any fixed integer shift \(t\):
\[
  \sum_{j=0}^{k}\alpha_jy_{n+j+t}
  =h\sum_{j=0}^{k}\beta_jf_{n+j+t}.
\]
Its associated shifted operator is
\[
  L_t[y(x);h]
  =\sum_{j=0}^{k}\alpha_jy\bigl(x+(j+t)h\bigr)
  -h\sum_{j=0}^{k}\beta_jy'\bigl(x+(j+t)h\bigr).
\]
Since \(L_t[y(x);h]=L[y(x+th);h]\), if \(L\) has order \(p\), then
\[
\begin{split}
  L_t[y(x);h]
  &=c_{p+1}h^{p+1}y^{(p+1)}(x+th)+O(h^{p+2})\\
  &=c_{p+1}h^{p+1}y^{(p+1)}(x)+O(h^{p+2}).
\end{split}
\]

\begin{definition}[Error constant]
Provided \(\sigma(1)=\sum_{j=0}^{k}\beta_j\neq0\), the error constant of
an order-\(p\) method is
\[
  C=\frac{c_{p+1}}{\sigma(1)}.
\]
\end{definition}

\begin{definition}[Consistency]
A linear multistep method is consistent if it has order at least one.
Equivalently,
\[
  c_0=c_1=0,
  \qquad\text{or}\qquad
  \rho(1)=0,\quad \rho'(1)=\sigma(1).
\]
\end{definition}

\subsection{Local Truncation Error}

Suppose \(y\) is an exact solution of
\[
  y'=f(x,y).
\]
The normalized local truncation error of the multistep method is
\[
  T_n
  =\frac{
    \displaystyle\sum_{j=0}^{k}
    \left[\alpha_jy(x_{n+j})-h\beta_jy'(x_{n+j})\right]
  }{
    \displaystyle h\sum_{j=0}^{k}\beta_j
  }.
\]
Taylor expansion about \(x_n\) yields
\[
  T_n
  =\frac{1}{\sigma(1)}
   \left[
     \frac{c_0}{h}y(x_n)+c_1y'(x_n)
     +c_2hy''(x_n)+\cdots
   \right].
\]
Thus consistency is equivalent to \(c_0=c_1=0\), and then
\(T_n\to0\) as \(h\to0\) along every sufficiently smooth solution
curve.

\begin{definition}[Consistency with the differential equation]
The numerical scheme is consistent with the differential equation if,
for every \(\varepsilon>0\), there exists \(h(\varepsilon)>0\) such that
\[
  |T_n|<\varepsilon
  \qquad\text{whenever}\qquad 0<h<h(\varepsilon),
\]
for any \(k+1\) consecutive points on a sufficiently smooth solution
curve of the initial-value problem.
\end{definition}

\begin{definition}[Order of accuracy]
The method has order \(p\) if \(p\) is the largest integer for which,
on every sufficiently smooth solution curve,
\[
  |T_n|\leq Kh^p
\]
for all sufficiently small \(h\), with \(K\) independent of \(h\).
\end{definition}

Consequently, the method has order \(p\) precisely when
\[
  c_0=c_1=\cdots=c_p=0,
  \qquad c_{p+1}\neq0.
\]
In that case,
\[
  T_n
  =\frac{c_{p+1}}{\sigma(1)}
   h^p y^{(p+1)}(x_n)+O(h^{p+1}),
\]
and \(c_{p+1}/\sigma(1)\) is the error constant.

\section{Difference Equations}

\begin{definition}
A difference equation is an equation involving differences, or shifted
values, of a sequence \(y_n\).  Its order is the difference between the
largest and smallest subscripts of \(y\) occurring in the equation.
\end{definition}

\begin{example}
The equation
\[
  \Delta^2y_n-2\Delta y_n+3y_n=n^2+n
\]
is equivalent to
\[
  (y_{n+2}-2y_{n+1}+y_n)
  -2(y_{n+1}-y_n)+3y_n=n^2+n,
\]
and hence to the second-order constant-coefficient equation
\[
  y_{n+2}-4y_{n+1}+6y_n=n^2+n.
\]
\end{example}

A linear difference equation of order \(k\) with constant coefficients
has the form
\[
  a_0y_{n+k}+a_1y_{n+k-1}+\cdots+a_{k-1}y_{n+1}+a_ky_n=f(n),
\]
where the coefficients \(a_0,\ldots,a_k\) are independent of \(n\).
It is homogeneous when \(f(n)=0\).

\subsection{Shift-Operator Form}

With \(Ey_n=y_{n+1}\), one has \(E^ry_n=y_{n+r}\), and the difference
equation can be written
\[
  (a_0E^k+a_1E^{k-1}+\cdots+a_k)y_n=f(n).
\]
Defining
\[
  \phi(E)=a_0E^k+a_1E^{k-1}+\cdots+a_k,
\]
its simplest form is
\[
  \phi(E)y_n=f(n).
\]

For the homogeneous equation \(\phi(E)y_n=0\), try a solution
\(y_n=\xi^n\).  Substitution gives
\[
  \xi^n\bigl(a_0\xi^k+a_1\xi^{k-1}+\cdots+a_k\bigr)=0,
\]
so \(\xi\) must satisfy the auxiliary, or characteristic, equation
\[
  \phi(\xi)
  =a_0\xi^k+a_1\xi^{k-1}+\cdots+a_k=0.
\]

\section{Linear Difference Equations}

\subsection{Homogeneous equations with constant coefficients}

Consider the homogeneous linear difference equation
\[
  a_0y_{n+k}+a_1y_{n+k-1}+\cdots+a_ky_n=0,
  \qquad a_0\ne 0.
\]
With the shift operator \(E\), defined by \(Ey_n=y_{n+1}\), this may be
written as
\[
  \phi(E)y_n=0,
  \qquad
  \phi(\xi)=a_0\xi^k+a_1\xi^{k-1}+\cdots+a_k.
\]
The equation \(\phi(\xi)=0\) is called the auxiliary, or characteristic,
equation.

If the roots \(\xi_1,\ldots,\xi_k\) are distinct, then
\[
  y_n=c_1\xi_1^n+c_2\xi_2^n+\cdots+c_k\xi_k^n,
\]
where \(c_1,\ldots,c_k\) are arbitrary constants.

\subsubsection{Repeated roots}

Suppose first that \(\xi_1\) is a double root.  Then
\[
  (E-\xi_1)^2y_n=0.
\]
Put \(y_n=\xi_1^n z_n\).  Since
\[
  y_{n+1}=\xi_1^{n+1}z_{n+1},
  \qquad
  y_{n+2}=\xi_1^{n+2}z_{n+2},
\]
we obtain
\[
  \xi_1^{n+2}(z_{n+2}-2z_{n+1}+z_n)=0.
\]
Thus
\[
  \Delta^2z_n=0,
  \qquad \Delta=E-1,
\]
and hence \(z_n=c_1+c_2n\).  The contribution associated with the
double root is therefore
\[
  y_n=(c_1+c_2n)\xi_1^n.
\]

More generally, if \(\xi_1\) has multiplicity \(r\), then its contribution
to the complementary solution is
\[
  \bigl(c_1+c_2n+\cdots+c_rn^{r-1}\bigr)\xi_1^n.
\]
The complete solution is obtained by adding the corresponding terms for
all the other roots.

\subsubsection{Complex-conjugate roots}

Suppose a real-coefficient equation has the conjugate roots
\[
  \xi_{1,2}=a\pm ib
  =r(\cos\theta\pm i\sin\theta),
  \qquad
  r=\sqrt{a^2+b^2},
\]
where \(\cos\theta=a/r\) and \(\sin\theta=b/r\).  Since
\[
  \xi_{1,2}^n=r^n\bigl(\cos(n\theta)\pm i\sin(n\theta)\bigr),
\]
the associated real-valued part of the complementary solution is
\[
  y_n=r^n\bigl(A\cos(n\theta)+B\sin(n\theta)\bigr),
\]
where \(A\) and \(B\) are arbitrary real constants.

\begin{example}
Solve
\[
  y_{n+2}+3y_{n+1}+2y_n=0,
  \qquad y_0=2,\quad y_1=1.
\]
\end{example}

\begin{proof}[Solution]
The operator equation is
\[
  (E^2+3E+2)y_n=(E+1)(E+2)y_n=0.
\]
The roots are \(-1\) and \(-2\), so
\[
  y_n=c_1(-1)^n+c_2(-2)^n.
\]
The initial conditions give
\[
  c_1+c_2=2,
  \qquad
  -c_1-2c_2=1.
\]
Therefore \(c_2=-3\), \(c_1=5\), and
\[
  \boxed{y_n=5(-1)^n-3(-2)^n.}
\]
\end{proof}

\begin{example}
Solve
\[
  u_{n+2}-2u_{n+1}+4u_n=0.
\]
\end{example}

\begin{proof}[Solution]
The characteristic equation
\[
  \xi^2-2\xi+4=0
\]
has roots
\[
  \xi=1\pm i\sqrt{3}
  =2\left(\cos\frac{\pi}{3}\pm i\sin\frac{\pi}{3}\right).
\]
Consequently,
\[
  \boxed{
  u_n=2^n\left(
    A\cos\frac{n\pi}{3}+B\sin\frac{n\pi}{3}
  \right).}
\]
\end{proof}

\begin{example}
Solve
\[
  u_{n+1}-2u_n+2u_{n-1}=0.
\]
\end{example}

\begin{proof}[Solution]
The characteristic roots are \(1\pm i\), whose modulus is \(\sqrt{2}\)
and whose arguments are \(\pm\pi/4\).  Hence
\[
  \boxed{
  u_n=(\sqrt{2})^n\left(
    A\cos\frac{n\pi}{4}+B\sin\frac{n\pi}{4}
  \right).}
\]
\end{proof}

\subsection{Nonhomogeneous equations}

Consider
\[
  \phi(E)y_n=f(n).
\]
If \(u_n\) is the complementary solution satisfying
\(\phi(E)u_n=0\), and \(v_n\) is any particular solution satisfying
\(\phi(E)v_n=f(n)\), then the complete solution is
\[
  y_n=u_n+v_n.
\]
Formally, a particular solution is denoted by
\[
  v_n=\frac{1}{\phi(E)}f(n).
\]

\subsubsection{Exponential forcing}

For \(f(n)=\beta^n\),
\[
  \phi(E)\beta^n=\phi(\beta)\beta^n.
\]
Thus, provided \(\phi(\beta)\ne0\),
\[
  \boxed{
  \frac{1}{\phi(E)}\beta^n
  =\frac{\beta^n}{\phi(\beta)}.}
\]
When \(\phi(\beta)=0\), this direct rule fails and the resonant case must
be treated by the shifted-operator formula developed below.

\subsection{Factorial polynomials and operator inversion}

\begin{definition}
The factorial polynomial of degree \(p\) is
\[
  n^{(p)}=n(n-1)\cdots(n-p+1),
\]
with \(n^{(0)}=1\) and \(n^{(1)}=n\).
\end{definition}

The forward-difference operator satisfies
\[
  \Delta n^{(p)}=p\,n^{(p-1)}
\]
and, more generally,
\[
  \Delta^m n^{(p)}
  =p^{(m)}n^{(p-m)}
  =p(p-1)\cdots(p-m+1)n^{(p-m)}
  \qquad (m\le p).
\]
In particular,
\[
  \Delta^p n^{(p)}=p!,
  \qquad
  \Delta^m n^{(p)}=0\quad(m>p).
\]
Formal inverse differences give, for example,
\[
  \frac{1}{\Delta}1=n^{(1)},\qquad
  \frac{1}{\Delta^2}1=\frac{n^{(2)}}{2!},\qquad
  \frac{1}{\Delta^3}1=\frac{n^{(3)}}{3!},
\]
up to terms belonging to the complementary solution.

If \(f(n)\) is a polynomial of degree at most \(m\), express it in the
factorial basis:
\[
  f(n)=A_0n^{(m)}+A_1n^{(m-1)}+\cdots+A_mn^{(0)}.
\]
Since \(E=1+\Delta\), a particular solution of
\(\phi(E)y_n=f(n)\) is
\[
  v_n=
  \frac{1}{\phi(1+\Delta)}
  \left(
    A_0n^{(m)}+A_1n^{(m-1)}+\cdots+A_mn^{(0)}
  \right).
\]
One expands \(1/\phi(1+\Delta)\) in ascending powers of \(\Delta\);
the resulting series terminates when it acts on a factorial polynomial.

\subsubsection{The shift rule}

For any polynomial \(\phi\) and sequence \(v(n)\),
\[
  \boxed{
  \phi(E)\bigl[\beta^n v(n)\bigr]
  =\beta^n\phi(\beta E)v(n).}
\]
Indeed, for every nonnegative integer \(j\),
\[
  E^j\bigl[\beta^n v(n)\bigr]
  =\beta^{n+j}v(n+j)
  =\beta^n(\beta E)^jv(n),
\]
and linearity gives the stated identity.  Consequently,
\[
  \boxed{
  \frac{1}{\phi(E)}\bigl[\beta^n v(n)\bigr]
  =\beta^n\frac{1}{\phi(\beta E)}v(n).}
\]
This rule is particularly useful in the resonant case
\(\phi(\beta)=0\): set \(E=1+\Delta\), expand
\(1/\phi(\beta(1+\Delta))\) in powers of \(\Delta\), and apply it to
\(v(n)\).

\subsection{Worked nonhomogeneous equations}

\begin{example}
Solve
\[
  2u_{n+2}+5u_{n+1}+2u_n=4.
\]
\end{example}

\begin{proof}[Solution]
The complementary equation factors as
\[
  2E^2+5E+2=(E+2)(2E+1),
\]
so
\[
  u_n^{(c)}=c_1(-2)^n+c_2\left(-\frac12\right)^n.
\]
For the constant forcing,
\[
  u_n^{(p)}
  =\frac{1}{2E^2+5E+2}\,4
  =\frac{4}{2+5+2}
  =\frac49.
\]
Hence
\[
  \boxed{
  u_n=c_1(-2)^n+c_2\left(-\frac12\right)^n+\frac49.}
\]
\end{proof}

\begin{example}
Solve
\[
  y_{n+3}-2y_{n+2}-4y_{n+1}+8y_n=3^n.
\]
\end{example}

\begin{proof}[Solution]
Here
\[
  E^3-2E^2-4E+8=(E-2)^2(E+2),
\]
so the complementary solution is
\[
  y_n^{(c)}=(c_1+c_2n)2^n+c_3(-2)^n.
\]
Since
\[
  \phi(3)=27-18-12+8=5,
\]
the particular solution is \(y_n^{(p)}=3^n/5\).  Therefore
\[
  \boxed{
  y_n=(c_1+c_2n)2^n+c_3(-2)^n+\frac{3^n}{5}.}
\]
\end{proof}

\begin{example}
Solve
\[
  (E^2-6E+8)y_n=3n+2.
\]
\end{example}

\begin{proof}[Solution]
The forcing may be written in factorial-polynomial form as
\[
  3n+2=3n^{(1)}+2n^{(0)}.
\]
The complementary solution is
\[
  y_n^{(c)}=c_1\,2^n+c_2\,4^n.
\]
A linear trial \(y_n^{(p)}=an+b\) gives
\[
  (E^2-6E+8)(an+b)=3an-4a+3b.
\]
Thus \(a=1\), \(b=2\), and
\[
  \boxed{y_n=c_1\,2^n+c_2\,4^n+n+2.}
\]
\end{proof}

\begin{example}
Solve
\[
  y_{n+3}-3y_{n+2}+3y_{n+1}-y_n=1.
\]
\end{example}

\begin{proof}[Solution]
The equation is
\[
  (E-1)^3y_n=\Delta^3y_n=1.
\]
Its complementary solution is
\[
  y_n^{(c)}=c_1+c_2n+c_3n^2.
\]
Since \(\Delta^3n^{(3)}=3!\), a particular solution is
\[
  y_n^{(p)}=\frac{n^{(3)}}{3!}
  =\frac{n(n-1)(n-2)}{6}.
\]
Hence
\[
  \boxed{
  y_n=c_1+c_2n+c_3n^2+\frac{n(n-1)(n-2)}{6}.}
\]
\end{proof}

\begin{example}
Solve
\[
  y_{n+2}-2y_{n+1}+y_n=3n+5.
\]
\end{example}

\begin{proof}[Solution]
Since the operator is \((E-1)^2=\Delta^2\), the complementary solution
is \(c_1+c_2n\).  Also,
\[
  3n+5=3n^{(1)}+5n^{(0)},
\]
and therefore
\[
  y_n^{(p)}
  =\frac{1}{\Delta^2}
    \bigl(3n^{(1)}+5n^{(0)}\bigr)
  =\frac{3}{3!}n^{(3)}+\frac{5}{2!}n^{(2)}.
\]
Thus
\[
  \boxed{
  y_n=c_1+c_2n
  +\frac12n(n-1)(n-2)
  +\frac52n(n-1).}
\]
\end{proof}

\section{Order of a Linear Multistep Method}

Consider the \(k\)-step linear multistep method
\[
  \sum_{j=0}^k\alpha_jy_{n+j}
  =h\sum_{j=0}^k\beta_jf_{n+j},
\]
with characteristic polynomials
\[
  \rho(z)=\sum_{j=0}^k\alpha_jz^j,
  \qquad
  \sigma(z)=\sum_{j=0}^k\beta_jz^j.
\]
For an exact solution \(y\), and assuming \(\sigma(1)\ne0\), define the
normalised local truncation error by
\[
  T_n=
  \frac{
    \displaystyle
    \sum_{j=0}^k
    \left[
      \alpha_jy(x_{n+j})-h\beta_jy'(x_{n+j})
    \right]
  }{
    \displaystyle h\sum_{j=0}^k\beta_j
  }
  =
  \frac{
    \displaystyle
    \sum_{j=0}^k
    \left[
      \alpha_jy(x_{n+j})-h\beta_jy'(x_{n+j})
    \right]
  }{h\sigma(1)}.
\]

Choose a branch of \(\log z\) in a neighbourhood of \(z=1\), and define
\[
  \Psi(z)=\frac{\rho(z)}{\log z}-\sigma(z)
  =\frac{\rho(z)-\sigma(z)\log z}{\log z}.
\]
The apparent singularity at \(z=1\) is removable whenever
\(\rho(1)=0\).

\begin{theorem}[Order characterisation]
Let \(p\) be a positive integer.  The linear multistep method has order
\(p\) if and only if \(\Psi\) has a zero of multiplicity \(p\) at
\(z=1\).
\end{theorem}

\begin{proof}
Suppose first that the method has order \(p\).  Its local truncation
error has the form
\[
  T_n=
  \frac{c_{p+1}}{\sigma(1)}
  h^p y^{(p+1)}(x_n)+O(h^{p+1}),
  \qquad c_{p+1}\ne0.
\]
Apply the method to the test problem
\[
  y'=y,\qquad y(0)=1,
\]
whose solution is \(y(x)=e^x\).  Since \(x_n=nh\),
\[
  \sum_{j=0}^k\alpha_jy(x_{n+j})
  =e^{nh}\rho(e^h),
\]
and
\[
  \sum_{j=0}^k\beta_jy'(x_{n+j})
  =e^{nh}\sigma(e^h).
\]
Consequently,
\[
  T_n
  =\frac{e^{nh}}{h\sigma(1)}
    \left[\rho(e^h)-h\sigma(e^h)\right]
  =e^{nh}\left(
    \frac{c_{p+1}}{\sigma(1)}h^p+O(h^{p+1})
  \right).
\]
Thus the analytic function
\[
  F(h)=\frac{\rho(e^h)-h\sigma(e^h)}{h}
      =\Psi(e^h)
\]
has a zero of multiplicity \(p\) at \(h=0\).  Since \(z=e^h\) is
biholomorphic between neighbourhoods of \(h=0\) and \(z=1\),
\(\Psi(z)\) has a zero of multiplicity \(p\) at \(z=1\).

Conversely, suppose that \(\Psi(z)\) has a zero of multiplicity \(p\)
at \(z=1\).  Then
\[
  F(h)=\Psi(e^h)
  =\frac{\rho(e^h)-h\sigma(e^h)}{h}
\]
has a zero of multiplicity \(p\) at \(h=0\).  Equivalently,
\[
  g(h):=\rho(e^h)-h\sigma(e^h)
  =\sum_{j=0}^k(\alpha_j-h\beta_j)e^{jh}
\]
has a zero of multiplicity \(p+1\) at \(h=0\).  Therefore
\[
  g^{(q)}(0)=0\quad(0\le q\le p),
  \qquad
  g^{(p+1)}(0)\ne0.
\]
Direct differentiation gives
\[
  g^{(q)}(0)
  =\sum_{j=0}^k
    \left(j^q\alpha_j-qj^{q-1}\beta_j\right)
  =q!c_q,
\]
where
\[
  c_0=\sum_{j=0}^k\alpha_j
\]
and, for \(q\ge1\),
\[
  c_q=
  \frac{1}{q!}\sum_{j=0}^k j^q\alpha_j
  -\frac{1}{(q-1)!}\sum_{j=0}^k j^{q-1}\beta_j.
\]
It follows that
\[
  c_0=c_1=\cdots=c_p=0,
  \qquad c_{p+1}\ne0,
\]
which are precisely the order-\(p\) conditions for a linear multistep
method.
\end{proof}

\begin{theorem}[Construction of the second characteristic polynomial]
Let \(\rho\) be a polynomial of degree \(k\) such that
\[
  \rho(1)=0,\qquad \rho'(1)\ne0.
\]
There is a unique polynomial \(\sigma\) of degree at most \(k\) for
which the linear multistep method associated with \(\rho\) and
\(\sigma\) has order at least \(k+1\).  In particular, consistency gives
\[
  \sigma(1)=\rho'(1)\ne0.
\]
\end{theorem}

\begin{proof}
Because \(\rho(1)=0\), the function
\[
  R(z)=\frac{\rho(z)}{\log z}
\]
has a removable singularity at \(z=1\), with \(R(1)=\rho'(1)\).
There is a unique polynomial of degree at most \(k\) whose
Taylor expansion about \(z=1\) agrees with that of \(R\) through degree
\(k\), namely
\[
  \sigma(z)=
  \sum_{m=0}^{k}
  \frac{R^{(m)}(1)}{m!}(z-1)^m.
\]
Hence \(R(z)-\sigma(z)\) has a zero of multiplicity at least
\(k+1\) at \(z=1\).  The order characterisation then proves
the result.
\end{proof}
\end{document}
